% Options for packages loaded elsewhere
% Options for packages loaded elsewhere
\PassOptionsToPackage{unicode}{hyperref}
\PassOptionsToPackage{hyphens}{url}
\PassOptionsToPackage{dvipsnames,svgnames,x11names}{xcolor}
%
\documentclass[
  letterpaper,
  DIV=11,
  numbers=noendperiod]{scrreprt}
\usepackage{xcolor}
\usepackage{amsmath,amssymb}
\setcounter{secnumdepth}{5}
\usepackage{iftex}
\ifPDFTeX
  \usepackage[T1]{fontenc}
  \usepackage[utf8]{inputenc}
  \usepackage{textcomp} % provide euro and other symbols
\else % if luatex or xetex
  \usepackage{unicode-math} % this also loads fontspec
  \defaultfontfeatures{Scale=MatchLowercase}
  \defaultfontfeatures[\rmfamily]{Ligatures=TeX,Scale=1}
\fi
\usepackage{lmodern}
\ifPDFTeX\else
  % xetex/luatex font selection
\fi
% Use upquote if available, for straight quotes in verbatim environments
\IfFileExists{upquote.sty}{\usepackage{upquote}}{}
\IfFileExists{microtype.sty}{% use microtype if available
  \usepackage[]{microtype}
  \UseMicrotypeSet[protrusion]{basicmath} % disable protrusion for tt fonts
}{}
\makeatletter
\@ifundefined{KOMAClassName}{% if non-KOMA class
  \IfFileExists{parskip.sty}{%
    \usepackage{parskip}
  }{% else
    \setlength{\parindent}{0pt}
    \setlength{\parskip}{6pt plus 2pt minus 1pt}}
}{% if KOMA class
  \KOMAoptions{parskip=half}}
\makeatother
% Make \paragraph and \subparagraph free-standing
\makeatletter
\ifx\paragraph\undefined\else
  \let\oldparagraph\paragraph
  \renewcommand{\paragraph}{
    \@ifstar
      \xxxParagraphStar
      \xxxParagraphNoStar
  }
  \newcommand{\xxxParagraphStar}[1]{\oldparagraph*{#1}\mbox{}}
  \newcommand{\xxxParagraphNoStar}[1]{\oldparagraph{#1}\mbox{}}
\fi
\ifx\subparagraph\undefined\else
  \let\oldsubparagraph\subparagraph
  \renewcommand{\subparagraph}{
    \@ifstar
      \xxxSubParagraphStar
      \xxxSubParagraphNoStar
  }
  \newcommand{\xxxSubParagraphStar}[1]{\oldsubparagraph*{#1}\mbox{}}
  \newcommand{\xxxSubParagraphNoStar}[1]{\oldsubparagraph{#1}\mbox{}}
\fi
\makeatother


\usepackage{longtable,booktabs,array}
\newcounter{none} % for unnumbered tables
\usepackage{calc} % for calculating minipage widths
% Correct order of tables after \paragraph or \subparagraph
\usepackage{etoolbox}
\makeatletter
\patchcmd\longtable{\par}{\if@noskipsec\mbox{}\fi\par}{}{}
\makeatother
% Allow footnotes in longtable head/foot
\IfFileExists{footnotehyper.sty}{\usepackage{footnotehyper}}{\usepackage{footnote}}
\makesavenoteenv{longtable}
\usepackage{graphicx}
\makeatletter
\newsavebox\pandoc@box
\newcommand*\pandocbounded[1]{% scales image to fit in text height/width
  \sbox\pandoc@box{#1}%
  \Gscale@div\@tempa{\textheight}{\dimexpr\ht\pandoc@box+\dp\pandoc@box\relax}%
  \Gscale@div\@tempb{\linewidth}{\wd\pandoc@box}%
  \ifdim\@tempb\p@<\@tempa\p@\let\@tempa\@tempb\fi% select the smaller of both
  \ifdim\@tempa\p@<\p@\scalebox{\@tempa}{\usebox\pandoc@box}%
  \else\usebox{\pandoc@box}%
  \fi%
}
% Set default figure placement to htbp
\def\fps@figure{htbp}
\makeatother





\setlength{\emergencystretch}{3em} % prevent overfull lines

\providecommand{\tightlist}{%
  \setlength{\itemsep}{0pt}\setlength{\parskip}{0pt}}



 


\KOMAoption{captions}{tableheading}
\makeatletter
\@ifpackageloaded{tcolorbox}{}{\usepackage[skins,breakable]{tcolorbox}}
\@ifpackageloaded{fontawesome5}{}{\usepackage{fontawesome5}}
\definecolor{quarto-callout-color}{HTML}{909090}
\definecolor{quarto-callout-note-color}{HTML}{0758E5}
\definecolor{quarto-callout-important-color}{HTML}{CC1914}
\definecolor{quarto-callout-warning-color}{HTML}{EB9113}
\definecolor{quarto-callout-tip-color}{HTML}{00A047}
\definecolor{quarto-callout-caution-color}{HTML}{FC5300}
\definecolor{quarto-callout-color-frame}{HTML}{acacac}
\definecolor{quarto-callout-note-color-frame}{HTML}{4582ec}
\definecolor{quarto-callout-important-color-frame}{HTML}{d9534f}
\definecolor{quarto-callout-warning-color-frame}{HTML}{f0ad4e}
\definecolor{quarto-callout-tip-color-frame}{HTML}{02b875}
\definecolor{quarto-callout-caution-color-frame}{HTML}{fd7e14}
\makeatother
\makeatletter
\@ifpackageloaded{bookmark}{}{\usepackage{bookmark}}
\makeatother
\makeatletter
\@ifpackageloaded{caption}{}{\usepackage{caption}}
\AtBeginDocument{%
\ifdefined\contentsname
  \renewcommand*\contentsname{Table of contents}
\else
  \newcommand\contentsname{Table of contents}
\fi
\ifdefined\listfigurename
  \renewcommand*\listfigurename{List of Figures}
\else
  \newcommand\listfigurename{List of Figures}
\fi
\ifdefined\listtablename
  \renewcommand*\listtablename{List of Tables}
\else
  \newcommand\listtablename{List of Tables}
\fi
\ifdefined\figurename
  \renewcommand*\figurename{Figure}
\else
  \newcommand\figurename{Figure}
\fi
\ifdefined\tablename
  \renewcommand*\tablename{Table}
\else
  \newcommand\tablename{Table}
\fi
}
\@ifpackageloaded{float}{}{\usepackage{float}}
\floatstyle{ruled}
\@ifundefined{c@chapter}{\newfloat{codelisting}{h}{lop}}{\newfloat{codelisting}{h}{lop}[chapter]}
\floatname{codelisting}{Listing}
\newcommand*\listoflistings{\listof{codelisting}{List of Listings}}
\makeatother
\makeatletter
\makeatother
\makeatletter
\@ifpackageloaded{caption}{}{\usepackage{caption}}
\@ifpackageloaded{subcaption}{}{\usepackage{subcaption}}
\makeatother
\usepackage{bookmark}
\IfFileExists{xurl.sty}{\usepackage{xurl}}{} % add URL line breaks if available
\urlstyle{same}
\hypersetup{
  pdftitle={Solid Rocket Internal Ballistics},
  pdfauthor={Troy Altus},
  colorlinks=true,
  linkcolor={blue},
  filecolor={Maroon},
  citecolor={Blue},
  urlcolor={Blue},
  pdfcreator={LaTeX via pandoc}}


\title{Solid Rocket Internal Ballistics}
\usepackage{etoolbox}
\makeatletter
\providecommand{\subtitle}[1]{% add subtitle to \maketitle
  \apptocmd{\@title}{\par {\large #1 \par}}{}{}
}
\makeatother
\subtitle{Theory ⋅ Equations ⋅ Numerical Simulation}
\author{Troy Altus}
\date{2026-04-05}
\begin{document}
\maketitle

\renewcommand*\contentsname{Table of contents}
{
\setcounter{tocdepth}{2}
\tableofcontents
}

\bookmarksetup{startatroot}

\chapter{Solid Rocket Internal
Ballistics}\label{solid-rocket-internal-ballistics}

Theory ⋅ Equations ⋅ Numerical Simulation

\hfill\break

\section{Welcome}\label{welcome}

This book is a comprehensive guide to the \textbf{internal ballistics}
of solid propellant rocket motors. It covers the fundamental theory,
mathematical derivations, and practical numerical simulation using
Python.

\subsection{What This Book Covers}\label{what-this-book-covers}

\begin{itemize}
\tightlist
\item
  Propellant characteristics and burn rate laws
\item
  Grain geometry and burnback simulation
\item
  Chamber pressure dynamics and K/N ratio
\item
  Erosive burning, ignition transients, and nozzle performance
\item
  A complete 0D internal ballistics simulator
\end{itemize}

\subsection{How to Use This Book}\label{how-to-use-this-book}

\textbf{Start here:} - \href{chapters/01-intro.qmd}{Chapter 1:
Introduction} --- Motivation and overview -
\href{chapters/05-internal-ballistics.qmd}{Chapter 5: Internal
Ballistics} --- Core 0D model -
\href{chapters/10-full-simulation.qmd}{Chapter 10: Full Simulation} ---
Complete working model

Use the \textbf{left sidebar} to navigate through all chapters and
appendices.

The accompanying Python code (\texttt{code/}) will eventually provide a
reusable simulation library.

\begin{tcolorbox}[enhanced jigsaw, arc=.35mm, bottomrule=.15mm, bottomtitle=1mm, breakable, colback=white, colbacktitle=quarto-callout-tip-color!10!white, colframe=quarto-callout-tip-color-frame, coltitle=black, left=2mm, leftrule=.75mm, opacityback=0, opacitybacktitle=0.6, rightrule=.15mm, title=\textcolor{quarto-callout-tip-color}{\faLightbulb}\hspace{0.5em}{Tip}, titlerule=0mm, toprule=.15mm, toptitle=1mm]

\textbf{Tip}: Many chapters contain executable code examples. You can
run them directly in the HTML version using \texttt{quarto\ preview}.

\end{tcolorbox}

\begin{center}\rule{0.5\linewidth}{0.5pt}\end{center}

\textbf{Repository}:
https://github.com/altustd/solid-rocket-internal-ballistics

\bookmarksetup{startatroot}

\chapter*{Preface}\label{preface}

\markboth{Preface}{Preface}

Welcome to this ebook on solid rocket internal ballistics. This book is
intended for students, engineers, and hobbyists interested in
understanding the fundamental principles governing the performance of
solid propellant rocket motors. We will cover the key concepts of
burning rate, Kn ratio, grain design, and performance modeling, with a
focus on practical applications and computational methods. Whether you
are designing your first motor or looking to deepen your understanding
of internal ballistics, this book will provide a solid foundation and
guide you through the complexities of solid rocket motor behavior. Let's
get started!

\bookmarksetup{startatroot}

\chapter{Introduction to Solid Rocket Internal
Ballistics}\label{introduction-to-solid-rocket-internal-ballistics}

Solid rocket motors are widely used in missiles, launch vehicles, and
amateur rocketry due to their simplicity, reliability, and high
thrust-to-weight ratio.

\section{What is Internal
Ballistics?}\label{what-is-internal-ballistics}

Internal ballistics is the study of the processes occurring
\textbf{inside} the rocket motor from ignition until burnout. It
includes:

\begin{itemize}
\tightlist
\item
  Ignition and flame spreading
\item
  Propellant combustion
\item
  Gas generation and pressure buildup
\item
  Nozzle flow and thrust generation
\end{itemize}

\section{Why It Matters}\label{why-it-matters}

Accurate internal ballistics modeling allows engineers to: - Predict
pressure-time curves - Optimize grain geometry - Ensure structural
integrity - Maximize performance and safety

This book focuses on the theoretical foundations and practical
computational methods used in solid rocket motor design.

\bookmarksetup{startatroot}

\chapter{Solid Propellants}\label{solid-propellants}

\section{Introduction}\label{introduction}

Solid propellants are energetic materials that contain both fuel and
oxidizer in a solid matrix. They are the heart of solid rocket motors,
providing the chemical energy converted into thrust.

\section{Classification of Solid
Propellants}\label{classification-of-solid-propellants}

\subsection{Double-Base Propellants}\label{double-base-propellants}

\begin{itemize}
\tightlist
\item
  Composed primarily of nitrocellulose and nitroglycerin
\item
  Homogeneous (uniform composition)
\item
  Common in older tactical missiles
\end{itemize}

\subsection{Composite Propellants}\label{composite-propellants}

\begin{itemize}
\tightlist
\item
  Most widely used today
\item
  Heterogeneous: oxidizer crystals (usually AP --- ammonium perchlorate)
  suspended in a fuel binder (HTPB, PBAN, etc.)
\item
  Higher performance and better mechanical properties
\end{itemize}

\subsection{Composite Modified Double-Base
(CMDB)}\label{composite-modified-double-base-cmdb}

\begin{itemize}
\tightlist
\item
  Hybrid of the two above
\end{itemize}

\section{Key Ingredients}\label{key-ingredients}

\begin{itemize}
\tightlist
\item
  \textbf{Oxidizer}: Ammonium Perchlorate (AP), Ammonium Nitrate (AN)
\item
  \textbf{Fuel/Binder}: Hydroxyl-terminated polybutadiene (HTPB),
  Polybutadiene acrylonitrile (PBAN)
\item
  \textbf{Metal Fuel}: Aluminum powder (increases specific impulse and
  density)
\item
  \textbf{Burn Rate Modifiers}: Iron oxide, copper chromite
\item
  \textbf{Plasticizers and Curatives}
\end{itemize}

\section{Important Properties}\label{important-properties}

\begin{itemize}
\tightlist
\item
  Burning rate (typically 5--25 mm/s at 70 bar)
\item
  Pressure exponent (\texttt{n} in Saint-Robert's law: \(r = a P^n\))
\item
  Density
\item
  Mechanical properties (tensile strength, elongation)
\item
  Sensitivity (shock, friction, thermal)
\end{itemize}

\section{Burning Rate Law (Vieille's / Saint-Robert's
Law)}\label{burning-rate-law-vieilles-saint-roberts-law}

\[
r = a P^n
\]

Where: - \(r\) = linear burning rate (mm/s) - \(P\) = chamber pressure
(bar or MPa) - \(a\) = burn rate coefficient (pressure and temperature
dependent) - \(n\) = pressure exponent (usually 0.2 -- 0.7)

\section{Next Steps}\label{next-steps}

In the following chapters we will explore how propellant properties,
grain geometry, and nozzle design interact to determine the internal
ballistics behavior of a solid rocket motor.

\bookmarksetup{startatroot}

\chapter{Grain Design and Geometry}\label{grain-design-and-geometry}

\section{Introduction}\label{introduction-1}

Grain design is one of the most critical aspects of solid rocket motor
performance. The shape of the propellant grain controls how the burning
surface area evolves over time, which directly determines the thrust
profile of the motor.

\section{Key Concepts}\label{key-concepts}

\subsection{Web Thickness}\label{web-thickness}

The \textbf{web thickness} is the minimum distance the burning surface
must regress before the grain is fully consumed.

\subsection{\texorpdfstring{Burning Surface Area
\(A_b(t)\)}{Burning Surface Area A\_b(t)}}\label{burning-surface-area-a_bt}

The evolution of burning surface area is the core variable in grain
design:

\[
A_b(t) = f(\text{grain geometry}, \text{web distance burned})
\]

Grain geometries produce three main burning behaviors: -
\textbf{Neutral}: Roughly constant thrust - \textbf{Progressive}: Thrust
increases with time - \textbf{Regressive}: Thrust decreases with time

\section{Common Grain Geometries}\label{common-grain-geometries}

\begin{itemize}
\tightlist
\item
  Cylindrical (core-burning)
\item
  Star grains
\item
  Finocyl (fin-in-cylinder)
\item
  Wagon wheel / dendrite
\item
  Conical and tapered grains
\end{itemize}

\section{Important Design Parameters}\label{important-design-parameters}

\begin{itemize}
\tightlist
\item
  Kn ratio (\(K_n = A_b / A_t\))
\item
  Volumetric loading fraction
\item
  Length-to-diameter (L/D) ratio
\item
  Sliver fraction (unburned propellant at burnout)
\end{itemize}

\section{Link to Internal Ballistics}\label{link-to-internal-ballistics}

Grain geometry combined with burning rate laws determines the
pressure-time curve. A well-designed grain allows engineers to tailor
thrust profiles for specific missions.

The next chapter explores the detailed equations of internal ballistics.

\bookmarksetup{startatroot}

\chapter{Burn Surface Regression and Grain
Burnback}\label{burn-surface-regression-and-grain-burnback}

\bookmarksetup{startatroot}

\chapter{Burn Surface Regression and Grain Burnback
Simulation}\label{burn-surface-regression-and-grain-burnback-simulation}

\section{Why Burnback Simulation
Matters}\label{why-burnback-simulation-matters}

The propellant grain geometry determines how the burning surface evolves
over time. Accurate modeling of this \textbf{burnback} (surface
regression) is essential for predicting realistic chamber pressure-time
(\(P\)-\(t\)) and thrust-time curves.

Simple analytical models work well for basic shapes such as cylindrical
core burners or end-burners. However, modern rocket motors often use
complex geometries --- star, finocyl, wagon-wheel, dendrite, or slotted
grains --- that require numerical burnback simulation.

\section{Core Concept: Web Distance}\label{core-concept-web-distance}

As the propellant burns, the burning surface moves \textbf{inward},
perpendicular to itself, by a distance called the \textbf{web} (\(w\)).

The instantaneous burning surface area \(A_b\) is a function of the web
distance burned:

\[
A_b(w) = f(\text{initial grain geometry}, w)
\]

Burnout occurs when the web distance \(w\) reaches the maximum web
thickness of the grain.

\section{Burnback Modeling
Approaches}\label{burnback-modeling-approaches}

\subsection{1. Analytical Methods (Fast but
Limited)}\label{analytical-methods-fast-but-limited}

Suitable for simple geometries:

\begin{itemize}
\tightlist
\item
  Internal-burning cylinder
\item
  External-burning cylinder
\item
  End-burner (cigarette grain)
\item
  Simple conical or tapered grains
\end{itemize}

These use closed-form equations but break down quickly for complex
ports.

\subsection{2. Numerical / Geometric Burnback (General
Purpose)}\label{numerical-geometric-burnback-general-purpose}

For arbitrary grain shapes, the standard techniques are:

\begin{itemize}
\tightlist
\item
  \textbf{Minimum Distance Function (MDF)} method
\item
  \textbf{Level-set} methods
\item
  Polygon offsetting (using libraries such as \texttt{shapely} in
  Python)
\item
  Voxel-based or grid-based regression
\item
  CAD-based parametric slicing
\end{itemize}

For most amateur, experimental, and small-motor work, a \textbf{2D
cross-sectional + axial slicing} approach (Ballistic Element Analysis)
offers an excellent balance of accuracy and speed.

\section{How Ballistic Element Analysis
Works}\label{how-ballistic-element-analysis-works}

\begin{enumerate}
\def\labelenumi{\arabic{enumi}.}
\tightlist
\item
  Slice the grain into many thin axial elements of thickness
  \(\Delta x\).
\item
  For each slice, compute the 2-D burning perimeter \(P_b(w)\) as a
  function of web distance.
\item
  Sum the contributions:
\end{enumerate}

\[
A_b(w) = \sum_{i=1}^{N} P_{b,i}(w) \cdot \Delta x
\]

\begin{enumerate}
\def\labelenumi{\arabic{enumi}.}
\setcounter{enumi}{3}
\tightlist
\item
  Simultaneously track the free chamber volume \(V_c(w)\) by subtracting
  burned propellant volume.
\end{enumerate}

This produces high-resolution tables of \(A_b(w)\) and \(V_c(w)\) that
are fed directly into the 0D lumped-parameter ballistic solver.

\section{Pseudocode: Basic Burnback + Ballistic
Coupling}\label{pseudocode-basic-burnback-ballistic-coupling}

\begin{figure}[H]

{\centering \pandocbounded{\includegraphics[keepaspectratio]{chapters/04-burnback-simulation_files/figure-pdf/cell-2-output-1.pdf}}

}

\caption{Example of burning surface area evolution during grain
burnback}

\end{figure}%

\bookmarksetup{startatroot}

\chapter{Internal Ballistics}\label{internal-ballistics}

0D Modeling of Chamber Pressure and Thrust-Time Evolution

\hfill\break

\section{Introduction}\label{introduction-2}

Having defined propellant properties and grain geometry regression
(burnback) in previous chapters, we now combine them to predict the time
evolution of chamber pressure, mass flow rate, and thrust.

This chapter introduces the \textbf{zero-dimensional (0D) internal
ballistics model}, the standard engineering approach for performance
prediction.

\section{Mass Balance in the Combustion
Chamber}\label{mass-balance-in-the-combustion-chamber}

The fundamental equation is conservation of mass inside the chamber:

\[
\frac{dm_c}{dt} = \dot{m}_\text{gen} - \dot{m}_\text{out}
\]

where: - ( m\_c ): mass of combustion products in the chamber (kg) - (
\dot{m}\_\text{gen} = \rho\emph{p A\_b(t) r(P\_c) ): mass generation
rate from burning propellant - ( \dot{m}}\text{out} ): mass expulsion
rate through the nozzle

\section{Burn Rate Law (Vieille's
Law)}\label{burn-rate-law-vieilles-law}

The linear burn rate is given by:

\[
r(P_c) = a P_c^n
\]

where ( a ) is the burn rate coefficient and ( n ) is the pressure
exponent (typically 0.2--0.7 for stable motors).

\section{Nozzle Mass Flow Rate}\label{nozzle-mass-flow-rate}

For a choked nozzle the outflow rate is:

\[
\dot{m}_\text{out} = \frac{P_c A_t}{c^*}
\]

where ( c\^{}* ) is the characteristic velocity.

\section{Governing Differential
Equation}\label{governing-differential-equation}

Combining the above leads to the core pressure evolution equation:

\[
\frac{V_c}{c^{*2}} \frac{dP_c}{dt} = \rho_p A_b(w) r(P_c) - \frac{P_c A_t}{c^*}
\]

We solve this coupled system numerically for chamber pressure ( P\_c(t)
) and web distance ( w(t) ), using ( A\_b(w) ) from the burnback module.

\section{Summary}\label{summary}

This chapter establishes the fundamental ODEs that link grain regression
to motor performance. The next chapter explores the K/N ratio and
equilibrium pressure in detail.

\begin{tcolorbox}[enhanced jigsaw, arc=.35mm, bottomrule=.15mm, bottomtitle=1mm, breakable, colback=white, colbacktitle=quarto-callout-note-color!10!white, colframe=quarto-callout-note-color-frame, coltitle=black, left=2mm, leftrule=.75mm, opacityback=0, opacitybacktitle=0.6, rightrule=.15mm, title=\textcolor{quarto-callout-note-color}{\faInfo}\hspace{0.5em}{Note}, titlerule=0mm, toprule=.15mm, toptitle=1mm]

Executable examples are available in the \texttt{code/} directory.

\end{tcolorbox}

\bookmarksetup{startatroot}

\chapter{Burning Rate and K/N Ratio}\label{burning-rate-and-kn-ratio}

Equilibrium Pressure, Stability, and Motor Ballistic Behavior

\hfill\break

\section{Introduction}\label{introduction-3}

In Chapter 5 we derived the governing differential equation for chamber
pressure evolution using mass conservation. At any instant, the chamber
pressure is determined by the balance between \textbf{mass generation}
from the burning propellant surface and \textbf{mass discharge} through
the nozzle.

When the time derivative of pressure is zero (( \frac{dP_c}{dt}
\approx 0 )), the motor operates at \textbf{equilibrium pressure}. This
condition dominates most of the burn duration in well-designed motors
and is controlled primarily by two quantities:

\begin{itemize}
\tightlist
\item
  The propellant \textbf{burn rate law}
\item
  The \textbf{K/N ratio} (( K\_n = A\_b / A\_t )), the ratio of burning
  surface area to nozzle throat area
\end{itemize}

This chapter explores these relationships in detail, their implications
for motor stability, pressure sensitivity, and how grain geometry (via
burnback) shapes the pressure-time and thrust-time curves.

\section{Vieille's Law (Saint-Robert's
Law)}\label{vieilles-law-saint-roberts-law}

The linear burn rate ( r ) of most solid propellants follows the
empirical power-law relationship known as \textbf{Vieille's law}:

\[
r(P_c) = a \, P_c^n
\]

where: - ( r ): linear burn (regression) rate (m/s or mm/s) - ( a ):
burn rate coefficient (units depend on pressure units, e.g., m/s·Paⁿ) -
( n ): dimensionless pressure exponent (typically 0.2--0.7 for stable
rocket propellants) - ( P\_c ): chamber pressure (usually expressed in
MPa or Pa)

The coefficient ( a ) is moderately sensitive to the initial propellant
temperature ( T\_i ):

\[
a(T_i) = a_\text{ref} \left[1 + \sigma_p (T_i - T_\text{ref})\right]
\]

where ( \sigma\_p ) is the temperature sensitivity of burn rate
(typically 0.1--0.3 \%/K).

A propellant is considered \textbf{ballistically stable} only when ( n
\textless{} 1 ). If ( n \geq 1 ), small pressure perturbations can cause
runaway pressure increase, often leading to motor explosion.

\section{Definition of K/N Ratio}\label{definition-of-kn-ratio}

The \textbf{K/N ratio} (commonly written ( K\_n ) or simply ( K )) is
defined as:

\[
K_n = \frac{A_b}{A_t}
\]

\begin{itemize}
\tightlist
\item
  ( A\_b ): instantaneous burning surface area (m²), obtained from the
  burnback model (Chapter 4)
\item
  ( A\_t ): nozzle throat area (m²), a fixed design parameter
\end{itemize}

( K\_n ) is one of the most important design parameters in solid rocket
motors. It directly controls the operating pressure level. Typical
values for amateur and experimental motors range from 100 to 400,
depending on propellant and desired pressure.

\section{Equilibrium Pressure
Derivation}\label{equilibrium-pressure-derivation}

At steady state, mass generation equals mass outflow:

\[
\dot{m}_\text{gen} = \dot{m}_\text{out}
\]

\[
\rho_p \, A_b \, r(P_c) = \frac{P_c A_t}{c^*}
\]

Substitute Vieille's law ( r = a P\_c\^{}n ):

\[
\rho_p \, A_b \, a \, P_c^n = \frac{P_c A_t}{c^*}
\]

Rearrange:

\[
\rho_p \, a \, c^* \, \left( \frac{A_b}{A_t} \right) = P_c^{1-n}
\]

\[
P_{eq} = \left( \rho_p \, a \, c^* \, K_n \right)^{\frac{1}{1-n}}
\]

This is the \textbf{equilibrium chamber pressure}. It shows that:

\begin{itemize}
\tightlist
\item
  Pressure scales strongly with ( K\_n ) because of the exponent (
  1/(1-n) ) (often 1.25--3.3).
\item
  Doubling ( K\_n ) typically increases pressure by a factor of
  \textasciitilde2.5--3 or more, depending on ( n ).
\item
  Higher ( n ) makes pressure much more sensitive to changes in ( K\_n
  ).
\end{itemize}

\section{K/N Curve and Pressure-Time
Profile}\label{kn-curve-and-pressure-time-profile}

As the grain burns, ( A\_b ) evolves according to the burnback geometry.
Plotting ( K\_n(w) ) versus web distance (or time) produces the
\textbf{K/N curve}.

The shape of this curve determines the ballistic classification of the
motor:

\begin{itemize}
\tightlist
\item
  \textbf{Neutral} --- ( K\_n ) roughly constant → nearly constant
  pressure and thrust (desired for many applications)
\item
  \textbf{Progressive} --- ( K\_n ) increases with time → rising
  pressure/thrust
\item
  \textbf{Regressive} --- ( K\_n ) decreases with time → falling
  pressure/thrust
\item
  \textbf{Progressive-regressive} (e.g., many star or finocyl grains)
  --- initial rise followed by a drop
\end{itemize}

Because ( P\_\{eq\} \propto K\_n\^{}\{1/(1-n)\} ), even modest changes
in ( K\_n ) produce noticeable changes in pressure. Grain designers
therefore carefully tailor the geometry so the K/N curve yields the
desired thrust profile.

\section{Stability Considerations}\label{stability-considerations}

For stable operation, two conditions must be satisfied:

\begin{enumerate}
\def\labelenumi{\arabic{enumi}.}
\tightlist
\item
  ( n \textless{} 1 ) (inherent propellant property)
\item
  The operating point must lie on the stable intersection of the mass
  generation curve (( \dot{m}\emph{\text{gen} \propto P\^{}n )) and the
  mass discharge curve (( \dot{m}}\text{out} \propto P ))
\end{enumerate}

If ( n \textgreater{} 1 ), the mass generation curve becomes steeper
than the discharge curve, and any small positive pressure perturbation
causes further increase --- unstable.

Temperature sensitivity also plays a role. The temperature sensitivity
of chamber pressure ( \pi\_K ) is related to the burn rate temperature
sensitivity ( \sigma\_p ) by:

\[
\pi_K = \frac{\sigma_p}{1 - n}
\]

This explains why motors with higher ( n ) are more sensitive to ambient
temperature changes.

\section{Practical Design
Implications}\label{practical-design-implications}

\begin{itemize}
\tightlist
\item
  \textbf{Choosing ( n )}: Lower ( n ) gives more stable, less
  temperature-sensitive motors but can make it harder to achieve high
  thrust.
\item
  \textbf{Choosing ( K\_n )}: Higher ( K\_n ) gives higher operating
  pressure (and usually higher thrust and specific impulse) but
  increases structural loads and erosive burning risk.
\item
  \textbf{Grain geometry}: The burnback module (Chapter 4) must provide
  accurate ( A\_b(w) ) so the resulting ( K\_n(t) ) produces the target
  pressure curve.
\item
  \textbf{Nozzle sizing}: Once the desired average pressure and
  propellant are chosen, ( A\_t ) is sized to achieve the target ( K\_n
  ).
\end{itemize}

In the accompanying simulation code, we compute equilibrium pressure
analytically when possible, but solve the full transient ODE numerically
to capture ignition, tail-off, and deviations caused by changing ( A\_b
).

\section{Summary}\label{summary-1}

The interplay between the burn rate law ( r = a P\^{}n ) and the K/N
ratio ( K\_n = A\_b / A\_t ) determines the equilibrium chamber pressure
via:

\[
P_{eq} = \left( \rho_p \, a \, c^* \, K_n \right)^{\frac{1}{1-n}}
\]

Understanding and controlling the K/N curve through grain geometry is
one of the most powerful tools available to the motor designer for
shaping thrust-time performance.

In the next chapter we will introduce erosive burning corrections, which
become important when high port velocities alter the effective burn rate
beyond the simple pressure dependence.

\section{References}\label{references}

\begin{itemize}
\tightlist
\item
  Sutton, G. P., \& Biblarz, O. (2017). \emph{Rocket Propulsion
  Elements}.
\item
  Greatrix, D. R. (2012). \emph{Powered Flight -- The Engineering of
  Aerospace Propulsion}.
\item
  Nakka, R. (various). Solid Rocket Motor Design Resources.
\item
  Kubota, N. (2004). Principles of Solid Rocket Motor Design.
\end{itemize}

\begin{tcolorbox}[enhanced jigsaw, arc=.35mm, bottomrule=.15mm, bottomtitle=1mm, breakable, colback=white, colbacktitle=quarto-callout-note-color!10!white, colframe=quarto-callout-note-color-frame, coltitle=black, left=2mm, leftrule=.75mm, opacityback=0, opacitybacktitle=0.6, rightrule=.15mm, title=\textcolor{quarto-callout-note-color}{\faInfo}\hspace{0.5em}{Note}, titlerule=0mm, toprule=.15mm, toptitle=1mm]

Executable examples showing K/N curve generation from different grain
geometries and the resulting pressure-time traces are available in the
\texttt{examples/} and \texttt{code/} directories.

\end{tcolorbox}

\bookmarksetup{startatroot}

\chapter{Erosive Burning}\label{erosive-burning}

Effects of High Port Velocity on Burning Rate

\hfill\break

\section{Introduction}\label{introduction-4}

In the standard 0D internal ballistics model, the burning rate depends
only on chamber pressure via Vieille's law ( r = a P\_c\^{}n ).

However, when hot combustion gases flow at high velocity parallel to the
burning surface, they can significantly increase the local burn rate.
This is called \textbf{erosive burning}. It is most severe near the aft
end of the grain where mass flux is highest.

\section{Lenoir-Robillard Model}\label{lenoir-robillard-model}

The total burn rate is modeled as:

\[
r = r_0 + r_e
\]

where ( r\_0 = a P\_c\^{}n ) and the erosive term is:

\[
r_e = \frac{\alpha G^{0.8}}{D^{0.2}} \exp\left( -\beta \frac{\rho_p r}{G} \right)
\]

\begin{itemize}
\tightlist
\item
  ( G = \dot{m} / A\_p ): local mass flux (kg/m²·s)
\item
  ( D ): hydraulic diameter of the port (m)
\item
  ( \alpha, \beta ): propellant-specific constants
\end{itemize}

\section{Effect on Mass Generation}\label{effect-on-mass-generation}

When erosive burning is active, the mass generation rate becomes:

\[
\dot{m}_\text{gen} = \rho_p A_b(w) \, r(P_c, G)
\]

This term is used in the chamber pressure differential equation from
Chapter 5.

\section{Summary}\label{summary-2}

Erosive burning is an important correction for long or high-performance
motors. It is usually added after the basic 0D model is working.

The next chapter discusses nozzle flow and thrust.

\begin{tcolorbox}[enhanced jigsaw, arc=.35mm, bottomrule=.15mm, bottomtitle=1mm, breakable, colback=white, colbacktitle=quarto-callout-note-color!10!white, colframe=quarto-callout-note-color-frame, coltitle=black, left=2mm, leftrule=.75mm, opacityback=0, opacitybacktitle=0.6, rightrule=.15mm, title=\textcolor{quarto-callout-note-color}{\faInfo}\hspace{0.5em}{Note}, titlerule=0mm, toprule=.15mm, toptitle=1mm]

Future Python code in \texttt{code/simulator.py} will support an
\texttt{erosive=True} option.

\end{tcolorbox}

\bookmarksetup{startatroot}

\chapter{Nozzle Flow and Thrust}\label{nozzle-flow-and-thrust}

Characteristic Velocity, Thrust Coefficient, and Motor Performance

\hfill\break

\section{Introduction}\label{introduction-5}

Chapters 5--7 focused on the combustion chamber: mass generation from
the propellant grain, burn rate laws, equilibrium pressure via the K/N
ratio, and corrections such as erosive burning.

Now we turn to the \textbf{nozzle}, which converts the high-pressure,
high-temperature combustion gases into high-velocity exhaust, producing
thrust.

This chapter covers the key nozzle relations used in internal
ballistics:

\begin{itemize}
\tightlist
\item
  Characteristic velocity ( c\^{}* ) (a measure of propellant combustion
  efficiency)
\item
  Thrust coefficient ( C\_F ) (a measure of nozzle expansion efficiency)
\item
  The overall thrust equation
\item
  Specific impulse ( I\_\{sp\} )
\end{itemize}

These quantities allow us to convert chamber pressure ( P\_c(t) ) and
mass flow rate directly into thrust ( F(t) ).

\section{Characteristic Velocity ( c\^{}*
)}\label{characteristic-velocity-c}

The \textbf{characteristic velocity} ( c\^{}* ) is a propellant property
that quantifies the efficiency of the combustion process and the
conversion of thermal energy into directed flow at the throat. It is
independent of the nozzle expansion ratio.

It is defined as:

\[
c^* = \frac{P_c A_t}{\dot{m}}
\]

where: - ( P\_c ): chamber pressure (Pa) - ( A\_t ): nozzle throat area
(m²) - ( \dot{m} ): mass flow rate through the nozzle (kg/s)

For an ideal gas with isentropic choked flow at the throat, the
theoretical ( c\^{}* ) is:

\[
c^* = \sqrt{ \frac{R T_c}{\gamma} } \left( \frac{\gamma + 1}{2} \right)^{\frac{\gamma + 1}{2(\gamma - 1)}} \frac{1}{\Gamma(\gamma)}
\]

or more commonly expressed using the gamma function form:

\[
c^* = \frac{\sqrt{\gamma R T_c}}{\gamma} \left( \frac{\gamma + 1}{2} \right)^{\frac{\gamma + 1}{2(\gamma - 1)}}
\]

where: - ( \gamma ): ratio of specific heats of the combustion products
- ( R ): specific gas constant (J/kg·K) - ( T\_c ): chamber (stagnation)
temperature (K)

In practice, ( c\^{}* ) values for solid propellants are obtained from
thermochemical equilibrium codes (CEA, RPA, etc.) and stored in the
propellant database. Typical values range from 1400--1800 m/s for
amateur propellants (e.g., KNSB, APCP).

\section{Thrust Equation}\label{thrust-equation}

The instantaneous thrust produced by the motor is:

\[
F = \dot{m} \, v_e + (P_e - P_a) A_e
\]

where: - ( v\_e ): effective exhaust velocity at the nozzle exit - (
P\_e ): exit plane pressure - ( P\_a ): ambient pressure - ( A\_e ):
nozzle exit area

For choked flow, ( \dot{m} = \frac{P_c A_t}{c^*} ). Substituting and
rearranging yields the very useful form:

\[
F = C_F \, P_c \, A_t
\]

where ( C\_F ) is the \textbf{thrust coefficient}.

\section{Thrust Coefficient ( C\_F )}\label{thrust-coefficient-c_f}

The thrust coefficient is a dimensionless parameter that captures the
efficiency of the nozzle expansion process. For an ideally expanded
nozzle (( P\_e = P\_a )) it depends primarily on ( \gamma ) and the
expansion ratio ( \epsilon = A\_e / A\_t ).

The ideal thrust coefficient is:

\[
C_F = \sqrt{ \frac{2 \gamma^2}{\gamma - 1} \left( \frac{2}{\gamma + 1} \right)^{\frac{\gamma + 1}{\gamma - 1}} \left[ 1 - \left( \frac{P_e}{P_c} \right)^{\frac{\gamma - 1}{\gamma}} \right] } + \left( \frac{P_e}{P_c} - \frac{P_a}{P_c} \right) \epsilon
\]

In vacuum (( P\_a = 0 )) and for optimum expansion (( P\_e = P\_a )),
this simplifies significantly.

In the 0D ballistics model, we usually assume a fixed ( C\_F )
(calculated from design expansion ratio and ( \gamma )) or compute it as
a weak function of ( P\_c ).

\section{Specific Impulse ( I\_\{sp\} )}\label{specific-impulse-i_sp}

Specific impulse is the most important overall performance metric:

\[
I_{sp} = \frac{F}{\dot{m} \, g_0} = \frac{C_F \, c^*}{g_0}
\]

where ( g\_0 = 9.80665 ) m/s².

Thus:

\begin{itemize}
\tightlist
\item
  ( c\^{}* ) reflects combustion chamber performance
\item
  ( C\_F ) reflects nozzle performance
\item
  Their product determines delivered ( I\_\{sp\} )
\end{itemize}

Typical delivered ( I\_\{sp\} ) for amateur solid motors ranges from
120--250 seconds depending on propellant and nozzle design.

\section{Integration into the Ballistics
Model}\label{integration-into-the-ballistics-model}

In the simulator (Chapter 5), once ( P\_c(t) ) is known from the ODE
solver:

\begin{enumerate}
\def\labelenumi{\arabic{enumi}.}
\tightlist
\item
  Compute mass flow: ( \dot{m}(t) = \frac{P_c(t) A_t}{c^*} )
\item
  Compute thrust: ( F(t) = C\_F , P\_c(t) , A\_t )
\item
  (Optional) Include ambient pressure correction if ( P\_e \neq P\_a )
\end{enumerate}

During tail-off (when web distance reaches maximum), ( A\_b ) drops
rapidly and thrust decays.

The total impulse is obtained by integrating thrust over burn time:

\[
I_{total} = \int_0^{t_b} F(t) \, dt
\]

\section{Key Design Considerations}\label{key-design-considerations}

\begin{itemize}
\tightlist
\item
  \textbf{Expansion ratio ( \epsilon )}: Higher ( \epsilon ) increases (
  C\_F ) (and ( I\_\{sp\} )) in vacuum but can cause over-expansion and
  flow separation at sea level.
\item
  \textbf{Throat erosion}: Real nozzles experience some throat
  enlargement, which decreases ( P\_c ) and alters performance. This can
  be modeled with a time-varying ( A\_t ).
\item
  \textbf{Nozzle efficiency losses}: Divergence loss, friction,
  two-phase flow (for metallized propellants), and non-ideal gas
  behavior reduce delivered ( C\_F ) and ( c\^{}* ) below theoretical
  values.
\item
  \textbf{Ambient pressure effects}: Motors fired at altitude deliver
  higher ( I\_\{sp\} ) than at sea level.
\end{itemize}

\section{Summary}\label{summary-3}

The nozzle transforms chamber conditions into thrust via two fundamental
parameters:

\begin{itemize}
\tightlist
\item
  ( c\^{}* ): propellant combustion figure of merit
\item
  ( C\_F ): nozzle expansion figure of merit
\end{itemize}

Together they give:

\[
F(t) = C_F \, P_c(t) \, A_t \qquad \text{and} \qquad I_{sp} = \frac{C_F \, c^*}{g_0}
\]

With ( P\_c(t) ) from the internal ballistics solver and ( A\_b(w) )
from the burnback model, we can now generate complete pressure-time and
thrust-time curves for any grain geometry and propellant combination.

In the next chapter we will address the \textbf{ignition transient} ---
the rapid pressurization phase from igniter activation until steady
burning is established.

\section{References}\label{references-1}

\begin{itemize}
\tightlist
\item
  Sutton, G. P., \& Biblarz, O. (2017). \emph{Rocket Propulsion
  Elements} (9th ed.).
\item
  Greatrix, D. R. (2012). \emph{Powered Flight -- The Engineering of
  Aerospace Propulsion}.
\item
  Nakka, R. Solid Rocket Motor Theory pages (characteristic velocity and
  thrust coefficient sections).
\item
  NASA Glenn Research Center Rocket Propulsion Basics.
\end{itemize}

\begin{tcolorbox}[enhanced jigsaw, arc=.35mm, bottomrule=.15mm, bottomtitle=1mm, breakable, colback=white, colbacktitle=quarto-callout-note-color!10!white, colframe=quarto-callout-note-color-frame, coltitle=black, left=2mm, leftrule=.75mm, opacityback=0, opacitybacktitle=0.6, rightrule=.15mm, title=\textcolor{quarto-callout-note-color}{\faInfo}\hspace{0.5em}{Note}, titlerule=0mm, toprule=.15mm, toptitle=1mm]

The Python implementation in \texttt{code/thermo.py} and
\texttt{code/simulator.py} includes functions to compute ( c\^{}* ), (
C\_F ), and thrust from propellant properties (( \gamma ), ( T\_c ), (
\epsilon )) and instantaneous chamber pressure. Examples will
demonstrate sea-level vs.~vacuum performance.

\end{tcolorbox}

\bookmarksetup{startatroot}

\chapter{Ignition Transient}\label{ignition-transient}

Pressurization Phase, Ignition Delay, Flame Spreading, and Initial
Pressure Spike

\hfill\break

\section{Introduction}\label{introduction-6}

All previous chapters assumed that steady or quasi-steady burning is
already established. In reality, every solid rocket motor begins with a
distinct \textbf{ignition transient} --- the short but critical phase
between igniter activation and the establishment of stable chamber
pressure.

During this phase (typically 0.01--0.5 seconds for amateur motors,
longer for large motors), the chamber pressure rises rapidly from
ambient to operating level, often overshooting the equilibrium pressure
before settling. This ``ignition spike'' can stress the case, grain, and
nozzle, and must be carefully controlled in design.

This chapter describes the physical processes, key stages, and
simplified modeling approaches for the ignition transient.

\section{Stages of the Ignition
Transient}\label{stages-of-the-ignition-transient}

The ignition transient is commonly divided into three overlapping
phases:

\begin{enumerate}
\def\labelenumi{\arabic{enumi}.}
\item
  \textbf{Ignition Induction (Delay) Phase}\\
  The igniter releases hot gases and/or particles into the chamber. Heat
  is transferred (primarily by convection and radiation) to the
  propellant surface. No significant propellant burning occurs yet.\\
  The propellant surface temperature rises until it reaches the
  \textbf{ignition temperature} ( T\_\{ig\} ) (typically 500--800 K
  depending on the propellant).\\
  This phase ends when the first portion of the grain ignites.
\item
  \textbf{Flame Spreading Phase}\\
  Once ignited at one location (usually near the igniter or head end),
  the flame spreads along the propellant surface.\\
  Flame spreading is driven by convective heating from the hot gases,
  enhanced by increasing chamber pressure and flow velocity.\\
  Erosive effects and local geometry strongly influence the spreading
  rate.
\item
  \textbf{Chamber Filling (Pressurization) Phase}\\
  As more surface area ignites, mass generation increases rapidly. The
  chamber fills with combustion products until the mass outflow through
  the nozzle balances generation.\\
  Pressure rises sharply, often overshooting the steady-state
  equilibrium value due to the sudden increase in burning surface before
  the flow stabilizes.
\end{enumerate}

After these phases, the motor transitions into the steady ballistic
phase described in Chapters 5--8.

\section{Igniter Modeling}\label{igniter-modeling}

The igniter is usually modeled as a time-dependent mass and energy
source:

\[
\dot{m}_{ig}(t), \quad T_{ig}, \quad \text{or total energy input}
\]

Simple models assume a constant or triangular igniter mass flow rate for
a short duration. More advanced models include particle impingement and
radiative effects.

The heat flux to the propellant surface during induction is dominated by
convection:

\[
\dot{q}'' = h (T_g - T_s)
\]

where ( h ) is the convective heat transfer coefficient (often using
correlations similar to those in erosive burning models), ( T\_g ) is
gas temperature, and ( T\_s ) is local surface temperature.

\section{Propellant Heating and Ignition
Criterion}\label{propellant-heating-and-ignition-criterion}

During the induction phase, the propellant is treated as inert. The
surface temperature rise is governed by the one-dimensional heat
conduction equation in the solid phase:

\[
\rho_p c_p \frac{\partial T}{\partial t} = k \frac{\partial^2 T}{\partial x^2}
\]

with boundary condition at the surface incorporating the heat flux from
the igniter gases.

A common practical ignition criterion is:

\begin{itemize}
\tightlist
\item
  The propellant surface reaches a critical ignition temperature (
  T\_\{ig\} ), or
\item
  The integrated energy input reaches a threshold value.
\end{itemize}

Once ( T\_s \geq T\_\{ig\} ), that surface element begins burning
according to Vieille's law (possibly with erosive correction).

\section{Simplified 0D Ignition Transient
Model}\label{simplified-0d-ignition-transient-model}

For engineering purposes in the 0D ballistics simulator, a
lumped-parameter approach can be used:

\begin{itemize}
\tightlist
\item
  Add igniter mass flow ( \dot{m}\_\{ig\}(t) ) to the mass balance
  equation.
\item
  Delay the start of propellant burning until a computed ignition delay
  time ( t\_\{ign\} ).
\item
  Ramp up the effective burning area ( A\_b(t) ) during flame spreading
  (e.g., using an exponential or linear ramp function based on grain
  geometry).
\end{itemize}

A basic ODE extension during ignition includes the igniter term:

\[
\frac{V_c}{c^{*2}} \frac{dP_c}{dt} = \rho_p A_b(t) r(P_c) + \dot{m}_{ig}(t) - \frac{P_c A_t}{c^*}
\]

where ( A\_b(t) ) starts from zero (or a small initial value) and grows
as flame spreads.

\section{Causes and Control of Ignition
Spike}\label{causes-and-control-of-ignition-spike}

The initial pressure spike occurs because:

\begin{itemize}
\tightlist
\item
  Flame spreading can ignite a large fraction of the surface almost
  simultaneously.
\item
  The burning surface at ignition may be higher than the steady-state
  value for some geometries.
\item
  Nozzle flow may not yet be fully choked or established.
\end{itemize}

Design strategies to reduce harmful spikes include:

\begin{itemize}
\tightlist
\item
  Optimized igniter mass flow and energy (not too strong, not too weak)
\item
  Grain geometry that promotes controlled flame spreading (e.g.,
  progressive port opening)
\item
  Nozzle throat closure (burst diaphragm) that opens at a designed
  pressure
\item
  Use of ignition aids or inhibitors on parts of the grain
\end{itemize}

\section{Numerical Considerations}\label{numerical-considerations}

Full-fidelity modeling of ignition transients often requires quasi-1D or
multi-dimensional CFD with coupled gas and solid phases, flame spreading
algorithms, and detailed heat transfer. For the purposes of this book
and the accompanying Python simulator, we use a simplified 0D model with
an optional ignition delay and flame-spreading function. This provides
good engineering accuracy for preliminary design while remaining
computationally lightweight.

Future extensions can include position-dependent ignition along the
grain length for more accurate prediction in long grains.

\section{Summary}\label{summary-4}

The ignition transient consists of induction, flame spreading, and
chamber filling phases. It is driven by igniter output,
convective/radiative heat transfer, and propellant thermal response.
Proper modeling is essential to predict and limit the initial pressure
spike, ensuring structural integrity and repeatable motor performance.

With the ignition transient included, the full 0D internal ballistics
simulation (Chapters 5--9) can now generate complete pressure-time and
thrust-time curves from t = 0 (igniter firing) through burnout.

In the next chapter we will discuss \textbf{overall motor performance
metrics}, validation against experimental data, and sensitivity
analysis.

\section{References}\label{references-2}

\begin{itemize}
\tightlist
\item
  Peretz, A., et al.~(1973). The Starting Transient of Solid-Propellant
  Rocket Motors with High Internal Gas Velocities. NASA Technical
  Reports.
\item
  Sutton, G. P., \& Biblarz, O. (2017). \emph{Rocket Propulsion
  Elements}.
\item
  Greatrix, D. R. (2012). \emph{Powered Flight -- The Engineering of
  Aerospace Propulsion}.
\item
  Various AIAA papers on ignition transient modeling (e.g., flame
  spreading and igniter effects).
\end{itemize}

\begin{tcolorbox}[enhanced jigsaw, arc=.35mm, bottomrule=.15mm, bottomtitle=1mm, breakable, colback=white, colbacktitle=quarto-callout-note-color!10!white, colframe=quarto-callout-note-color-frame, coltitle=black, left=2mm, leftrule=.75mm, opacityback=0, opacitybacktitle=0.6, rightrule=.15mm, title=\textcolor{quarto-callout-note-color}{\faInfo}\hspace{0.5em}{Note}, titlerule=0mm, toprule=.15mm, toptitle=1mm]

The Python implementation in \texttt{code/simulator.py} includes an
optional \texttt{ignition\_model} parameter. It supports simple delay +
ramp functions for flame spreading as well as direct addition of igniter
mass flow. Examples will demonstrate ignition spike prediction and
mitigation strategies.

\end{tcolorbox}

\bookmarksetup{startatroot}

\chapter{}\label{section}

\cleardoublepage
\phantomsection
\addcontentsline{toc}{part}{Appendices}
\appendix

\chapter{}\label{section-1}

\chapter{}\label{section-2}

\chapter{}\label{section-3}

\chapter{}\label{section-4}




\end{document}
